\begin{progprob}
    One of the most used tools in the data scientist's toolkit is the
    \emph{histogram}. We often use it to get a sense of the distribution of a
    dataset. You'll remember them from DSC 10, where we used them quite a bit.
    In this problem, you'll be asked to design and implement an efficient
    algorithm that computes a histogram from sorted data.

    \textbf{Computing a histogram.} We saw how to compute a density histogram in DSC 10,
    but here's a quick reminder. Suppose you're given a list of $n$ points, along with
    $k$ bins of the form $[a, b)$, where $a$ is the start of the bin and $b$ is the end (we
    assume that $a$ is included in the bin, but $b$ is not). The density within a given
    bin is calculated using the formula:
    \[
        \frac{\text{\# of points in the bin}}{(\text{total \# of points}) \times (\text{bin width})}
    \]
    For example, suppose we are given 8 data points: $(1, 2, 3, 6, 7, 9, 10, 11)$ and
    3 bins: $[0, 4), [4, 8),$ and $[8, 12)$. Three of the eight points fall in the first bin,
    whose width is 4. The density in the first bin is therefore $\frac{3}{8 \times 4} = 0.09375$.
    On the other hand, only two points fall into the second bin, so the density there is
    $\frac{2}{8 \times 4} = 0.0625$. Three points fall into the last bin, so the density there
    is also $\frac{3}{8 \times 4} = 0.09375$.

    We can visualize these densities using a plot like the one shown below:

    \begin{center}
        \includegraphics[width=.4\textwidth]{\thisdir/include/histogram.png}
    \end{center}

    \textbf{The problem.} In a file named \texttt{histogram.py}, write a function called \python{histogram(points, bins)}
    which accepts two arguments:
    \begin{itemize}
        \item \python{points}: a \textbf{sorted} Python list of $n$ numbers, in order from smallest to largest
        \item \python{bins}: a Python list of $k$ tuples, each of the form $(a, b)$, where $a$ and $b$ are the start
            and end of a bin, respectively. You can assume that the bins are
            also sorted from left to right, and the endpoint of one bin is the
            start of the next. For example, the list \python{[(0, 4), (4, 8),
            (8, 12)]} denotes bins of $[0, 4), [4, 8),$ and $[8, 12)$. Each bin
            includes its start, but not its end.
    \end{itemize}

    The output of your function should be a Python list which contains the
    histogram's density within each bin. So if \python{bins} has $k$ elements,
    your function should return a list of $k$ numbers.

    For example, here's what your code should do on a simple input:

    \begin{minted}[autogobble]{python}
        >>> points = [1, 2, 3, 6, 7, 9, 10, 11]
        >>> bins = [(0, 4), (4, 8), (8, 12)]
        >>> histogram(points, bins)
        [0.09375, 0.0625, 0.09375]
    \end{minted}

    Your code will be tested not only for correctness, but also for efficiency.
    To get full credit, your function should take $\Theta(n)$ time, where $n$
    is the number of data points. It should take this time regardless of the
    number of bins! Note that there is a simple, inefficient solution that will
    take $\Theta(n \times k)$ time (for instance, if the number of bins is
    $n/10$, it would take $\Theta(n^2)$ time!). Your code should be much faster
    than that, and it should only need to loop through the data points
    \emph{once}. \emph{Hint}: make good use of the fact that the data and the
    bins are given to your function in \textbf{sorted} order.

    For this problem, you may not use \python{numpy} or any other imports. You can
    do this in pure Python! You may assume that the number of bins is between $1$
    and $n$, and that each of the data points falls into one of the bins.

    \textbf{About the autograder.} 
    
    This is the first ``Programming Problem'' of
    the quarter. \textbf{You must submit your code to Gradescope through Github} in a 
    separate assignment named ``Homework 03 - Programming Problem 01'' with 
    the same due date as the written homework. \textbf{The guide for setting up Github 
    and submitting your code through it is available on Campuswire.} 

    Programming problems are mostly autograded, which means that your code will
    need to pass several tests to receive full credit. Some of these tests are
    \emph{public}, meaning that you can see the result before the deadline.
    They usually check to make sure that your code has the correct filename and
    works on the simplest of examples. The private tests check your code on
    more complex examples, and usually will make sure that it has the correct
    time complexity. You can submit your code as many times as you'd like, and
    we suggest submitting it early and often! After uploading your code, be
    sure to stick around and see if it passes the public tests. If it doesn't,
    you might lose points.


    \begin{soln}
        First, let's look at the simplest solution. It will return the right
        densities, but it is inefficient. It looks like this:

        \inputminted{python}{\thisdir/include-solution/slow-histogram.py}

        That solution has a nested loop which, for each bin, loops through the
        data points, counting the number of points in the bin. This means that
        we loop through the data once for every bin, or a total of $n \times k$
        times. For example, if the number of bins is $n/10$, the time
        complexity of this code is $\Theta(n^2)$. This is not efficient, and we
        can do better!

        The above solution does not take advantage of the fact that the data
        and the bins are sorted. Here's a solution that does:

        \inputminted{python}{\thisdir/include-solution/histogram.py}

        This code also has a nested loop, but there's a key difference. For
        each bin, we start iterating through the data points, but stop when we
        see a data point past the end of the bin. The key observation is that,
        for the next bin, we can start where we left off previously, with the
        first data point that was past the end of the previous bin. This means
        that we only loop through the data once, and the time complexity is
        $\Theta(n)$, as required, no matter how many bins there are.
    \end{soln}

\end{progprob}
